
\section{Precisione}\label{sec:precisione}

Il progetto utilizza un unico tipo scalare, \code{scalar\_t}, definito in
\file{src/common/scalar.h}. Per default esso corrisponde a \code{double},
mentre la compilazione con \code{PREC=float} seleziona \code{float} tramite
la macro \code{USE\_FLOAT}.

Anche il tipo utilizzato nelle comunicazioni MPI viene derivato dalla stessa
scelta attraverso \code{SCALAR\_MPI\_TYPE}. La struttura dell'algoritmo e
l'interfaccia dei backend rimangono quindi invariate fra le build FP64 e
FP32, mentre cambia il tipo numerico sul quale vengono eseguite le
operazioni.

Passare da FP64 a FP32 dimezza la dimensione di ciascun elemento,
riducendo il volume di dati associato allo stesso numero di valori.
A parità di modello di traffico, ciò aumenta quindi l'intensità aritmetica
espressa in FLOP per byte.

\color{red}
La precisione modifica però anche la capacità di calcolo disponibile
sull'hardware, che può essere molto diversa fra FP32 e FP64. Il regime
prestazionale non può quindi essere dedotto dalla sola riduzione della
dimensione dei dati: deve essere determinato confrontando intensità
aritmetica, banda di memoria e throughput di calcolo della piattaforma.
Questa analisi viene svolta nel \cref{cap:roofline}.

La scelta della precisione influenza anche la validazione numerica.
Il confronto con il riferimento seriale utilizza tolleranze dipendenti dal
tipo scalare e dalla lunghezza della riduzione, anziché una soglia unica
per tutte le configurazioni. La definizione della tolleranza e la misura
dell'errore sono discusse nel \cref{cap:validazione}.
\color{black}

\subsection{Generazione deterministica}

I dati vengono costruiti mediante una funzione deterministica basata sul
mixer a 64 bit di SplitMix64. Nel progetto il mixer non viene utilizzato
come una sequenza pseudocasuale da avanzare elemento dopo elemento:
l'input viene invece ricavato direttamente dal seed, da un identificatore
della matrice e dall'indice globale dell'elemento.

In questo modo ogni valore può essere calcolato indipendentemente dagli
altri, senza mantenere uno stato condiviso del generatore.

Indicando con $\Sigma_A$ e $\Sigma_X$ due identificatori distinti dei
dataset, la chiave associata a ciascun elemento dipende dalla sua posizione
nella matrice globale:

\begin{equation}
A[i][j]
=
g\!\left(\text{seed},\Sigma_A,iN+j\right),
\qquad
X[j][c]
=
g\!\left(\text{seed},\Sigma_X,jk+c\right).
\label{eq:gen}
\end{equation}

Gli identificatori distinti impediscono che $A$ e $X$ utilizzino la stessa
sequenza di valori pur condividendo lo stesso seed.

Di conseguenza, il valore associato a una posizione globale non dipende
dal rank MPI che la genera, dalla forma della griglia di processi o
dall'ordine con cui gli elementi vengono prodotti.

L'uscita a 64 bit del mixer viene convertita in un valore floating point
nell'intervallo $[0,1)$ e successivamente riscalata in $[-1,1)$.
La generazione viene effettuata inizialmente in doppia precisione e il
risultato viene convertito in \code{scalar\_t} soltanto alla fine.

\subsection{Indipendenza dalla decomposizione}

La generazione basata sugli indici globali ha tre conseguenze utili.

\begin{itemize}
    \item Ogni processo può generare direttamente gli elementi appartenenti
    al proprio blocco locale, senza dover ricevere necessariamente una copia
    della matrice globale.

    \item Cambiando il numero di processi o la forma della griglia MPI,
    ciascuna posizione globale continua ad avere lo stesso valore. Le
    configurazioni parallele vengono quindi confrontate sullo stesso
    problema numerico.

    \item Il riferimento seriale può ricostruire autonomamente gli elementi
    necessari alla validazione a partire dallo stesso seed e dagli stessi
    indici globali.
\end{itemize}

I valori vengono riscalati nell'intervallo $[-1,1)$, producendo dati con
segni sia positivi sia negativi e distribuzione approssimativamente
centrata intorno allo zero.

Questa scelta evita casi di test costituiti esclusivamente da contributi
dello stesso segno e produce riduzioni numericamente più rappresentative
dei diversi ordini di accumulazione introdotti dalle implementazioni
parallele.

Non viene tuttavia assunta una legge deterministica per la crescita
dell'errore. La tolleranza utilizzata nella validazione viene definita
separatamente e verificata sperimentalmente, come discusso nel
\cref{cap:validazione}.

La generazione avviene inizialmente in \code{double}; nel caso FP32 il
valore viene successivamente convertito in \code{float}. Le build FP64 e
FP32 partono quindi dalla stessa sequenza deterministica di valori reali,
differendo per l'arrotondamento introdotto dalla rappresentazione scelta.